\documentclass[11pt,a4paper]{article}
\usepackage[T1]{fontenc}
\usepackage[utf8]{inputenc}
\usepackage{lmodern,amsmath,amssymb,amsthm,mathtools}
\usepackage[margin=25mm]{geometry}
\usepackage{microtype,booktabs,enumitem}
\usepackage[colorlinks=true,linkcolor=blue,urlcolor=blue]{hyperref}
\newtheorem{theorem}{Theorem}
\newcommand{\Z}{\mathbb Z}
\newcommand{\N}{\mathbb N}
\newcommand{\strip}{\mathrm{strip}}
\newcommand{\pos}[1]{[#1]_+}
\setlength{\parindent}{0pt}
\setlength{\parskip}{5pt}
\setlist{nosep,leftmargin=*}
\title{Deletion of the middle edge of an FG strip\\\large A theorem of \emph{Ice and mirrors}: formalization target}
\author{Cluster HMS}
\date{14 November 2026}
\begin{document}
\maketitle
\vspace{-1.5em}
This document specifies the theorem to be proved in Lean. Its mathematical source is
the strip-deletion theorem of \emph{Ice and mirrors}.
The source commit is
\begin{center}\texttt{7e310e01df64e38be583500fc81a8ef1a001a196}.\end{center}
The accompanying blueprint records the supplied proof. Neither document asserts that
the Lean formalization is complete.

\section{The strip and its integer matrix}
Fix integers
\[
 n\ge2,\qquad q=\lfloor n/2\rfloor,\qquad N\ge\max\{q,2\}.
\]
The strip has distinct polygon vertices $T_0,\ldots,T_N,B_0,\ldots,B_N$.
Its oriented triangles, with the displayed cyclic orders, are
\[
 \tau_{2i+1}=(T_{i+1},T_i,B_i),\qquad
 \tau_{2i+2}=(B_i,B_{i+1},T_{i+1})\qquad(0\le i<N).
\]
Put
\[
 e_{2i}=T_iB_i\quad(0\le i\le N),\qquad
 e_{2i+1}=T_{i+1}B_i\quad(0\le i<N),\qquad \ell=e_N.
\]
Thus $\tau_j$ lies between $e_{j-1}$ and $e_j$. The distinguished edge is
$e_N$ also when $N$ is odd; it is then a diagonal $T_{(N+1)/2}B_{(N-1)/2}$.

In an oriented triangle $(U_0,U_1,U_2)$ take the lattice points
\[
 (a,b,c)\in\Z_{\ge0}^3,\qquad a+b+c=n,
 \qquad (a,b,c)\notin\{(n,0,0),(0,n,0),(0,0,n)\}.
\]
The coordinates are the weights at $U_0,U_1,U_2$. Points on a shared edge
are identified when their weights at its two endpoints agree. For each
$a,b,c\ge0$ with $a+b+c=n-2$, add the three arrows
\begin{equation}\label{eq:cycle}
 (a+1,b+1,c)\longrightarrow(a,b+1,c+1)
 \longrightarrow(a+1,b,c+1)\longrightarrow(a+1,b+1,c).
\end{equation}
Sum the signed arrow multiplicities after gluing. Omit all lattice points
on $T_iT_{i+1}$ and $B_iB_{i+1}$, for $0\le i<N$, and their incident arrows.
These are the actual top and bottom boundary segments. The lattice points
on $e_0$ and $e_{2N}$ are retained.

Let $V_{\strip}$ be the resulting finite set. Its matrix is
\[
 B_{\strip}(u,v)=\#(u\longrightarrow v)-\#(v\longrightarrow u)\in\Z.
\]
In particular, opposite arrows cancel and $B_{\strip}(u,v)=-B_{\strip}(v,u)$.
Write $L(e_j)\subset V_{\strip}$ for the $n-1$ lattice points on $e_j$, and put
\[
 P=L(e_0)\cup L(e_{2N}),\qquad I=V_{\strip}\setminus P.
\]

\newpage
\section{Ambient matrices and deletion moves}
Let $V$ be any finite set and $\iota:V_{\strip}\hookrightarrow V$ an injection.
Extend the strip matrix by zero to $V$ and denote the extension by
$B_{\strip,\mathrm{ext}}$. The input matrix is
\begin{equation}\label{eq:ambient}
 B=B_{\strip,\mathrm{ext}}+E,\qquad
 E_{uv}=-E_{vu},\qquad
 E_{\iota(x),v}=0\quad(x\in I,\ v\in V).
\end{equation}
The block $E|_{V\setminus\iota(I)}$ is arbitrary. It includes arrows between
protected vertices, arrows from protected vertices to exterior vertices,
and arrows between exterior vertices. Consequently the full induced matrix
on $\iota(V_{\strip})$ need not equal $B_{\strip}$.

For a live vertex set $U\subseteq V$, an integer skew-symmetric matrix $A$
on $U$, and $k\in U$, define $\pos{x}=\max\{x,0\}$ and
\begin{equation}\label{eq:mutation}
 (\mu_k A)_{ij}=\begin{cases}
 -A_{ij},&i=k\text{ or }j=k,\\
 A_{ij}+\pos{A_{ik}}\pos{A_{kj}}-
             \pos{-A_{ik}}\pos{-A_{kj}},&i,j\ne k.
 \end{cases}
\end{equation}
A word $W=(k_1,\ldots,k_m)$ acts in the listed order, so
$\mu_W A=\mu_{k_m}\cdots\mu_{k_1}A$. The empty word acts as the identity.
Let $W^{\mathrm{rev}}=(k_m,\ldots,k_1)$.

A \emph{deletion move} at $s\in U$ consists of a word in $U\setminus\{s\}$
such that
\begin{equation}\label{eq:sign}
 \bigl(\forall v\in U,\ (\mu_W A)_{sv}\ge0\bigr)
 \quad\text{or}\quad
 \bigl(\forall v\in U,\ (\mu_W A)_{sv}\le0\bigr),
\end{equation}
followed by deletion of $s$ and the reversed word. Deletion means restriction
to $U\setminus\{s\}$. The output is
\begin{equation}\label{eq:undo}
 \mu_{W^{\mathrm{rev}}}
 \bigl((\mu_W A)|_{U\setminus\{s\}}\bigr)
 =A|_{U\setminus\{s\}}.
\end{equation}
This equality must be proved from mutation involutivity and compatibility
with restriction away from the mutation vertex, or obtained from verified
library results. The sign condition is tested against every live ambient
vertex, including the protected and exterior vertices.

\begin{theorem}\label{thm:target}
For every $n,N,V,\iota,E$ satisfying the conditions above, there is a finite
sequence of deletion moves for $B$ whose deleted set $Z$ satisfies
\[
 \iota(L(e_N))\subseteq Z\subseteq\iota(I).
\]
Every mutation vertex belongs to $\iota(I)$ and is present when mutated.
Each mutation word avoids its deletion vertex. The final live set and matrix are
\[
 V\setminus Z,\qquad B|_{V\setminus Z}.
\]
In particular, no vertex of $\iota(P)$ or $V\setminus\iota(V_{\strip})$
is mutated or deleted.
\end{theorem}

The quantifiers over the rank, strip length, exterior vertices, and permitted
matrix $E$ are unrestricted. Existence of suitable deletion moves is the
conclusion, not part of the input data.

\newpage
\section{The application to a triangulated surface}
Consider the mutable FG quiver of an oriented triangulated surface containing
the embedded strip above. Its top and bottom segments lie on the actual
surface boundary, whose lattice points have been omitted. Assume the strip
embeds without identifications of its distinct lattice vertices. The other
triangles meet the strip only along $e_0$, $e_{2N}$, or the omitted boundary.

The triangle sum for the ambient quiver splits as
\[
 B_S=B_{\strip,\mathrm{ext}}+
       \sum_{\tau\ \text{outside the strip}}B_{\tau,\mathrm{ext}}.
\]
No triangle outside the strip contains a retained lattice vertex in
$\iota(I)$. Therefore the second summand satisfies \eqref{eq:ambient}.
Theorem~\ref{thm:target} gives the theorem for this quiver.
The formal deliverable must include this application, or an explicit
combinatorial application theorem with these hypotheses verified from
the embedded triangle construction.

An internal grid or sum-type representation is permitted. Its vertices,
edge identifications, omitted boundary points, and matrix entries must be
identified with the glued construction in Section~1. In particular, a theorem
for an unspecified matrix with assumed local deletion properties does not
establish Theorem~\ref{thm:target}.

\section{Source correspondence and acceptance}
The proof-to-code map must give actual Lean declaration names for the
following mathematical objects and conclusions.
\begin{center}
\begin{tabular}{@{}p{.58\linewidth}p{.34\linewidth}@{}}
\toprule
Mathematical content & Source in \texttt{main.tex}\\
\midrule
Triangle lattice, cycles, gluing, boundary omission & lines 3134--3151\\
Strip, deletion moves, main theorem & lines 3368--3409\\
Edge depths and triangle subsets $C_p$ & lines 3498--3509\\
Triangle-interior deletion & lines 3575--3654\\
Two-triangle coordinates & lines 3767--3792\\
Common-edge deletion and omitted mutation & lines 3860--3980\\
Rounds and removal of $e_N$ & lines 4119--4161\\
\bottomrule
\end{tabular}
\end{center}
The map must also identify the ambient extension hypothesis, mutation
convention, live vertex sets, and the exact conclusion of
Theorem~\ref{thm:target}. If the proof uses different intermediate results,
record how they establish the same statement.

Acceptance requires pinned Lean and mathlib versions, a reproducible
successful build, and a fresh consumer file importing the delivered theorem.
Audit its transitive axioms with \texttt{\#print axioms}, together with the
local theorems used in its proof. The only allowed axioms are
\[
 \texttt{propext},\qquad \texttt{Classical.choice},\qquad \texttt{Quot.sound}.
\]
No \texttt{sorry}, \texttt{admit}, custom axiom, unsafe proof bypass, or
additional trust from native evaluation is permitted.

The exact-integer script \texttt{verify\_fg\_strips.py} passed its checks for
$2\le n\le8$ during preparation. This is debugging evidence only.
Its global construction fixes $N=n\mathbin{//}2$ and has no arbitrary
exterior attachment. For $n=2,3$ it gives $N=1$, outside the required $N\ge2$.
Finite checks cannot replace the quantified theorem. The final report must
separate completed Lean results from unresolved obligations and identify
any source-proof correction explicitly.
\end{document}
